\section{Threats to Validity}
Labels. The labels come only from reruns, so tests whose flakiness does not reproduce on our machine are labeled stable, and the number of runs matters: the 20-run round identified only 13 of the 30 flaky tests identified with 100 runs. Environmental disruptions can also create spurious flaky tests; we mitigated them by discarding contaminated rounds and repeating them in isolation. Test identity is reconstructed from display names and class hierarchies, and a small number of tests could not be resolved to source code (none of them flaky in the 100-run round). Internal. The 1:1 undersampling balances the test folds and makes precision and F1 optimistic (LOPO quantifies the difference); the same 30 positives appear in every draw, so the confidence intervals underestimate uncertainty and, with so few positives, the ordering among classifiers is not reliable; correlated tests (for example, three NLog tests that failed together) may inflate the positive class; and the two best models were chosen on the same data on which they are evaluated. Attribute tokens may encode project conventions, and removing them reduces the MCC of the best models by about 0.05-0.07. External and selection. Project selection was not random: it favored open-source libraries with unit tests and domains in which concurrency is common, so the 0.24\% rate and the results may not generalize to applications, UI tests, or proprietary code; moreover, flakiness is concentrated in nine repositories, and the results are specific to the .NET toolchain and machine used.